\textbf{Multimodality in imitation.} Implicit Behavioral Cloning (IBC)~\cite{florence2021implicit} represents the policy as an energy-based model and argues that this handles discontinuities and multimodality better than explicit regression; Ta et al.~\cite{ta2022conditional} analyse the gap between theory and practice of such conditional EBMs. Behavior Transformers~\cite{shafiullah2022behavior} cluster actions with k-means and predict a bin plus a continuous offset, cloning $k$ modes with one model. Mixture heads were used for goal-conditioned play imitation by Lynch et al.~\cite{lynch2019learning}. Urain et al.~\cite{urain2024deep} survey deep generative models for learning from multimodal demonstrations. Mazza et al.~\cite{mazza2026understanding} study how different parameterisations of generative BC policies represent multimodal expert behaviour and identify different bottlenecks per parameterisation; our controlled toy study is complementary in being purely empirical and including a matched mixture baseline.

\textbf{Diffusion policies.} Diffusion Policy~\cite{chi2023diffusion} applies DDPMs~\cite{ho2020denoising} (with improved noise schedules~\cite{nichol2021improved}) to visuomotor action generation, with receding-horizon action chunks. Follow-up work reduces its cost or changes the generative mechanism, e.g.\ IMLE Policy~\cite{rana2025imle} (implicit maximum likelihood, one-step sampling) and CARP~\cite{gong2024carp} (coarse-to-fine autoregressive prediction). We do not reproduce these methods; we reimplement the three simplest heads (regression, mixture, DDPM) and keep everything else equal.

\textbf{Gap.} We are not aware of a small study that sweeps start-state ambiguity, compares against a mixture head with matched training budget, and reports per-mode shares; the claims below are limited to this setting.
